% (C) 2026 Ralf Stephan, in collaboration with Claude Code.  CC0 1.0.
\chapter[Reachability, components, ranks]{Reachability, strong components and ranks}\label{ch:reach}

Most of the finite witnesses in the problems of Chapter~\ref{ch:overview} are
\emph{ranks}: functions \(\rho\colon V\to\N\) that never increase along an arc. This chapter
develops ranks from the strong-component structure of \cite[§1.5]{BJG09}. The main results
characterise when a rank exists that leaves at most one rank-keeping successor
(Theorem~\ref{thm:rank-complete}), show that it forces eventual periodicity
(Theorem~\ref{thm:evper}), and give the peel-rank criterion that locates all cycles
(Theorem~\ref{thm:peel}).

\section{Reachability}

Write \(u\leadsto v\) if there is a walk from \(u\) to \(v\), and \(u\leadsto_n v\) if there is
one of length exactly \(n\). Two vertices are in the same \emph{strong component}\index{strong component} if
\(u\leadsto v\) and \(v\leadsto u\). This is an equivalence relation. A vertex is \emph{cyclic}\index{cyclic}
if a walk of positive length leads from it back to itself, that is, if it has a successor
\(t\) with \(t\leadsto v\). Cyclic vertices form whole strong components.
Homomorphisms preserve \(\leadsto\) and send cyclic vertices to cyclic vertices.

\paragraph{Computing reachability.} Let \(A\) be finite. The \emph{ball} of radius \(k\)
around \(v\) is \(B_0=\{v\}\), \(B_{k+1}=B_k\cup\operatorname{out}(B_k)\). So
\(w\in B_k\) iff \(v\leadsto_n w\) for some \(n\le k\).

\begin{proposition}\label{prop:ball}
If \(B_{k+1}=B_k\), then \(B_k\) contains every vertex reachable from \(v\). If \(V\) is finite,
this happens for some \(k\le|V|\). Hence the vertices reachable from \(v\) are exactly those
in \(B_{|V|}\), and reachability is decidable.
\end{proposition}

\begin{proof}
A ball that equals its successor ball is closed under \(\operatorname{out}\), so it contains
every vertex reachable from \(v\). If \(B_{k+1}\neq B_k\) for all \(k\le|V|\), then
\(|B_k|\ge k+1\) for \(k\le|V|\) by induction, and \(|B_{|V|}|>|V|\) is impossible. Balls
increase with \(k\).
\end{proof}

Mathlib has no decision procedure for the reflexive--transitive closure of a relation, so
Proposition~\ref{prop:ball} is what makes \(\leadsto\) usable by \texttt{decide}.

\section{The two canonical ranks}

Let \(V\) and \(A\) be finite. The \emph{descendant count} and the \emph{predecessor count} are
\[
  \operatorname{desc}(v)=|\{u : v\leadsto u\}|,\qquad
  \operatorname{preds}(v)=|\{u : u\leadsto v\}| .
\]

\begin{proposition}[acyclic ordering, {\cite[Prop.~1.4.3, §1.5]{BJG09}}]\label{prop:desc}
Let \(s\to t\) be an arc.
\begin{enumerate}[label=(\alph*)]
\item \(\operatorname{desc}(t)\le\operatorname{desc}(s)\), with equality iff \(t\leadsto s\),
  i.e.\ iff \(s\) and \(t\) lie in one strong component.
\item \(\operatorname{preds}(s)\le\operatorname{preds}(t)\). If \(s\) and \(t\) lie in different
  components, then \(\{u: u\leadsto s\}\subsetneq\{u:u\leadsto t\}\).
\item \(\operatorname{desc}\) of the reverse graph is \(\operatorname{preds}\).
\end{enumerate}
\end{proposition}

\begin{proof}
(a) Descendants of \(t\) are descendants of \(s\). If the two finite sets have equal size, they
are equal, so \(s\), a descendant of itself, is a descendant of \(t\). Conversely, if
\(t\leadsto s\) the inclusion goes both ways. (b) is dual, with \(t\) a predecessor of itself
but not of \(s\). (c) holds because a walk in the reverse graph is a reversed walk.
\end{proof}

So \(\operatorname{desc}\) is non-increasing along arcs and strictly decreasing exactly across
components. In other words, it is a rank form of the acyclic ordering of the strong component
digraph. Equivalently, the condensation is acyclic. The predecessor count is the dual ordering,
increasing along arcs. One problem family uses one convention and another uses the other. This
book uses non-increasing ranks throughout.

\section{Ranks and branching}

\begin{definition}\label{def:rank}
A \emph{rank}\index{rank} of \(X\) is a \(\rho\colon V\to\N\) such that
\begin{enumerate}[label=(R\arabic*)]
\item \(\rho(t)\le\rho(s)\) for every arc \(s\to t\);
\item every vertex \(s\) has at most one successor \(t\) with \(\rho(t)=\rho(s)\).
\end{enumerate}
A \emph{labelled rank}\index{rank!labelled} replaces (R2) by: at most one arc \(s\xrightarrow{a}t\) with
\(\rho(t)=\rho(s)\). \(X\) has \emph{no branching}\index{no branching} if no vertex has two distinct successors in
its own strong component (labelled: two distinct arcs).
\end{definition}

\begin{theorem}\label{thm:rank-complete}
For finite \(V\) and \(A\), \(X\) has a rank iff it has no branching, and then
\(\operatorname{desc}\) is a rank. The labelled version holds likewise. For infinite \(V\), if
every vertex has finitely many descendants, the same holds with
\(\operatorname{desc}(v)=|\{u:v\leadsto u\}|\).
\end{theorem}

\begin{proof}
A rank is constant on strong components: along a walk it cannot increase, and a walk leads
back. So two successors of \(s\) in the component of \(s\) are both rank-keeping, and (R2)
makes them equal. Conversely, without branching, \(\operatorname{desc}\) satisfies (R1) by
Proposition~\ref{prop:desc}(a). By the same proposition, a successor that keeps
\(\operatorname{desc}\) lies in the component of \(s\), and there is at most one of those. For
infinite \(V\), use \(|\cdot|\) on finite sets in the same argument.
\end{proof}

No branching means that every strong component containing an arc is a single cycle
\cite[Def.~2.2.13]{LM95}. Chapter~\ref{ch:entropy} shows that this is exactly zero entropy.
Here we prove its first consequence.

\begin{theorem}[eventual periodicity]\label{thm:evper}\index{eventually periodic!walks}
Let \(\rho\) be a rank and \(s_0\to s_1\to\cdots\) an infinite walk such that only finitely many
vertices are reachable from \(s_0\). Then there are \(n_0\) and \(L>0\) with \(s_{m+L}=s_m\) for
all \(m\ge n_0\). If \(\rho\) is a labelled rank and \((s,\ell)\) an infinite path, the labels
\(\ell\) are eventually periodic with the same period.
\end{theorem}

\begin{proof}
The sequence \(\rho(s_n)\) does not increase, so it is constant from some \(n_1\) on. From then
on every step is rank-keeping, so by (R2) \(s_{m+1}\) is determined by \(s_m\). The walk stays
in a finite set, so \(s_p=s_q\) for some \(n_1\le p<q\). By induction \(s_{m+(q-p)}=s_m\) for
all \(m\ge p\). In the labelled case the arcs \(s_m\xrightarrow{\ell_m}s_{m+1}\) and
\(s_m=s_{m+L}\xrightarrow{\ell_{m+L}}s_{m+L+1}=s_{m+1}\) are both rank-keeping, so
\(\ell_{m+L}=\ell_m\).
\end{proof}

Theorem~\ref{thm:evper} is the engine behind the digit problem of Chapter~\ref{ch:overview}.
If the residues of \(\lfloor\xi a^n\rfloor\) modulo \(M\) are eventually units, they form an
infinite walk in \(D_{a,M}\). A rank of \(D_{a,M}\) makes the walk eventually periodic. For
\(a\le M\) the digit is determined by consecutive residues, so the digit word is eventually
periodic too, and \cite[Theorem~4]{Dub09} turns that into infinitely many composite terms.
The theorem is local: it uses only the part of the graph reachable from the walk.

\section{Peel ranks}

\begin{theorem}[peel rank]\label{thm:peel}
Let \(C\supseteq\{\text{cyclic vertices}\}\) and \(K\subseteq V\). Suppose
\(\rho\colon V\to\N\) satisfies \(\rho(t)\le\rho(s)\) for every arc \(s\to t\) inside \(C\), and
\(\rho(t)<\rho(s)\) for every such arc with \(s\notin K\). Then every cyclic vertex lies in
\(K\). The same holds with all inequalities and arcs reversed: \(\rho(s)\le\rho(t)\) inside
\(C\), strictly if \(t\notin K\).
\end{theorem}

\begin{proof}
Let \(r\notin K\) be cyclic, with a closed walk \(r\to t\leadsto r\). Every vertex of this
closed walk is cyclic, hence in \(C\). The first step drops the rank strictly and no later step
raises it, so \(\rho(r)\le\rho(t)<\rho(r)\). For the reversed statement, apply this to the
reverse graph, which has the same cyclic vertices.
\end{proof}

\begin{corollary}[rungs]\label{cor:rung}\index{rung}
Let \(f\colon X\to Y\) be a homomorphism and \(K'\) a set containing every cyclic vertex of
\(Y\). If \(\rho\) is a peel rank for \(K\) on \(C=f^{-1}(K')\), then every cyclic vertex of
\(X\) lies in \(K\).
\end{corollary}

\begin{proof}
Homomorphisms send cyclic vertices to cyclic vertices, so \(f^{-1}(K')\) contains every cyclic
vertex of \(X\). Then apply Theorem~\ref{thm:peel}.
\end{proof}

With \(f\) the reduction \(D_{a,M}\to D_{a,M'}\) (Example~\ref{ex:proj}),
Corollary~\ref{cor:rung} is one rung of a ladder of moduli \(M_1\mid M_2\mid\cdots\). Each rung
only needs to check the rank on the fibre over the previous core. This is how the base-\(7\)
case of the composite problem was checked.

\begin{proposition}[core form]\label{prop:coreform}
If every cyclic vertex lies in \(K\) and no vertex of \(K\) has two distinct successors in \(K\),
then \(X\) has no branching. For finite \(V\), \(A\) it therefore has a rank.
\end{proposition}

\begin{proof}
Two successors of \(v\) in the component of \(v\) are cyclic, as is \(v\), so all three lie
in \(K\).
\end{proof}

The converse fails. A rank can exist while the set of units is not deterministic: in base
\(4\) modulo \(6\) a rank exists, but the units do not satisfy the hypothesis of
Proposition~\ref{prop:coreform} with \(K\) the set of all units.

\begin{proposition}[rank depth]\label{prop:depth}\index{rank!depth}
Along a walk \(g_0\to\dots\to g_n\) on which \(\rho\) never increases, the number of steps with
\(\rho(g_{i+1})<\rho(g_i)\), plus \(\rho(g_n)\), is at most \(\rho(g_0)\).
\end{proposition}

\begin{proof}
Induction on \(n\): each strict drop lowers \(\rho\) by at least one.
\end{proof}

\section{Geodesic Kőnig}

\begin{theorem}\label{thm:geodesic-konig}\index{K\H{o}nig's lemma!geodesic}
Let \(A\) be finite. If infinitely many vertices are reachable from \(v\), then there is an
infinite walk \(v=s_0\to s_1\to\cdots\) whose vertices are pairwise distinct. In fact
\(\operatorname{dist}(v,s_n)=n\) for all \(n\).
\end{theorem}

\begin{proof}
Let \(\operatorname{dist}(v,w)\) be the length of a shortest walk from \(v\) to \(w\). Call an
arc \(u\to w\) \emph{geodesic} if \(u\) is reachable and
\(\operatorname{dist}(v,w)=\operatorname{dist}(v,u)+1\). Every reachable \(w\) is reached from
\(v\) by geodesic arcs in exactly \(\operatorname{dist}(v,w)\) steps. To see this, take the last
step \(u\to w\) of a shortest walk. Then \(\operatorname{dist}(v,u)\) is at most the length of
the walk minus one, and at least \(\operatorname{dist}(v,w)-1\), so the arc is geodesic, and we
induct. Balls are finite, so for each \(n\) some reachable \(w\) has
\(\operatorname{dist}(v,w)\ge n\). Hence the geodesic subgraph has walks of every length from
\(v\). Theorem~\ref{thm:konig} gives an infinite geodesic walk, along which the distance
increases by one at each step.
\end{proof}

This is the abstract half of the finiteness theorem for follower automata in
Chapter~\ref{ch:sofic}: an infinite automaton has an infinite path that never revisits a state.

\section*{Notes}
Strong components and their acyclic ordering are \cite[§1.5]{BJG09}. The two rank conventions
appear in the problems as descendant counts (the digit and divisor problems) and as predecessor
counts (the word-graph certificates of \cite{Oda26}). Theorem~\ref{thm:evper} for finite graphs
and Theorem~\ref{thm:peel} come from the digit problem. Theorem~\ref{thm:rank-complete} is
their converse, and it shows that a rank witness exists exactly when the graph has zero entropy
(Chapter~\ref{ch:entropy}).

\section*{Lean dictionary}
\begin{center}\small
\begin{tabular}{@{}>{\raggedright}p{0.3\textwidth}>{\raggedright\arraybackslash}p{0.64\textwidth}@{}}
\toprule
Text & Lean (\texttt{CarryGraph/Reach.lean}, \texttt{Rank.lean})\\
\midrule
reachability, components, cycles & \lean{Reach}, \lean{ReachIn}, \lean{SameSCC},
  \lean{sameSCC\_equivalence}, \lean{Cyc}, \lean{cyc\_iff}\\
homomorphisms & \lean{Hom.reach\_map}, \lean{Hom.cyc\_map}\\
Proposition~\ref{prop:ball} & \lean{ball}, \lean{mem\_ball}, \lean{exists\_ball\_stable},
  \lean{reachSet}, \lean{mem\_reachSet}, \lean{decReach}\\
Proposition~\ref{prop:desc} & \lean{desc}, \lean{preds}, \lean{desc\_eq\_iff},
  \lean{desc\_lt\_of\_not\_reach}, \lean{preds\_ssubset}, \lean{desc\_reverse}\\
Definition~\ref{def:rank} & \lean{IsRank}, \lean{IsLRank}, \lean{NoBranch}, \lean{LNoBranch}\\
Theorem~\ref{thm:rank-complete} & \lean{IsRank.noBranch}, \lean{noBranch\_isRank\_desc},
  \lean{exists\_isRank\_iff\_noBranch}, \lean{noBranch\_isRank\_ncard}\\
Theorem~\ref{thm:evper} & \lean{IsRank.evPer}, \lean{IsLRank.evPer}\\
Theorem~\ref{thm:peel} & \lean{cyc\_mem\_of\_peelRank}, \lean{cyc\_mem\_of\_peelRank\_rev}\\
Corollary~\ref{cor:rung} & \lean{CarryBridge.coreClosed\_of\_ladder'}\\
Proposition~\ref{prop:coreform} & \lean{noBranch\_of\_cyc\_det}\\
Proposition~\ref{prop:depth} & \lean{drops\_add\_le}\\
Theorem~\ref{thm:geodesic-konig} & \lean{exists\_injective\_walk\_of\_infinite\_reach}\\
\bottomrule
\end{tabular}
\end{center}
